\newpage
\phantomsection
\label{prf:injective-restriction-local-inverse}
\LRAProofFor{prop:injective-restriction-local-inverse}

\begin{remark*}[Return]
\hyperref[prop:injective-restriction-local-inverse]{Return to Proposition}
\end{remark*}

\begin{proposition*}[Injective Restriction Gives a Local Inverse]
Let $f:A\to B$ and $C\subseteq A$. If $f|_C:C\to B$ is injective, then
$f|_C:C\to f(C)$ is bijective and therefore has an inverse $f(C)\to C$.
\end{proposition*}

\begin{proof}[Professional Standard Proof]
\LRAProofBodyStart
The map $f|_C:C\to f(C)$ is surjective by
\hyperref[prop:restricted-map-onto-image]{the restricted-image proposition}.
It is injective by the hypothesis. Hence it is bijective, so the inverse
function exists.
\end{proof}

\begin{proof}[Detailed Learning Proof]
\LRAProofBodyStart
The restriction removes all inputs outside $C$. If no two remaining inputs
collide, and if the codomain is exactly the set of values those inputs produce,
then the restricted map matches each value in $f(C)$ with exactly one input in
$C$.
\end{proof}
\begin{remark*}[Proof structure]
\end{remark*}



\begin{dependencies}
\begin{itemize}
  \item \hyperref[prop:injective-restriction-local-inverse]{Proof target}
\end{itemize}
\end{dependencies}

\clearpage

